\section{Boundary certificate: what the reduction does and does not prove}
\label{sec:boundary-certificate}

This section records the logical endpoint of the argument.  It separates
three statements which have different observables, different natural
scales, and different quantifiers: the row-centered primitive-shell
residual, one prescribed prime-pair shift, and a finite bundle of candidate
shifts.  None may be
substituted for another without an additional theorem.

\subsection{The shell and the endpoint live on different scales}

For fixed $h\ne0$, put
\[
 \mathcal P_h(X):=\sum_n\Lambda(n)\Lambda(n+h)W(n/X).
\]
The fixed-shift upper-bound sieve gives
\begin{equation}
 \mathcal P_h(X)\ll_{h,W}X,
 \label{eq:endpoint-upper-scale}
\end{equation}
with the prime-power terms absorbed in the same bound; see, for example,
\citet{IwaniecKowalski2004}.  The Hardy--Littlewood prediction is the much
sharper asymptotic
\begin{equation}
 \mathcal P_h(X)=\mathfrak S(h)XI_0(W)+o(X),
 \label{eq:hl-endpoint-boundary}
\end{equation}
for admissible even $h$ \citep{HardyLittlewood1923}.  Thus the endpoint
problem has natural scale $X$.

By contrast, the complete row-centered primitive-shell residual proved
above has the
form
\begin{equation}
 \cR_h(X)
 =\left(\frac{\varphi(|h|)}{|h|}-\beta_h\right)
   XI_0(W)\log X+O_{h,W}(X).
 \label{eq:shell-versus-endpoint-scale}
\end{equation}
When $\delta_hI_0(W)\ne0$, its scale is $X\log X$.  This extra logarithm is factor-coordinate
multiplicity, not an extra logarithm in the endpoint correlation.
Moreover, $\cR_h$ is a signed, centered sum over primitive
factorizations, whereas $\mathcal P_h$ is a one-dimensional endpoint
observable, nonnegative only when $W\geq0$.  No positivity-preserving comparison between the two quantities
has been proved.  In particular, dividing
\eqref{eq:shell-versus-endpoint-scale} by $\log X$ does not yield
\eqref{eq:hl-endpoint-boundary}; even the positive coefficient $1/6$ when
$h=2$ is not a twin-prime lower bound.  This is the same label-preservation
boundary isolated in TPC-10 and TPC-14
\citep{WangTPC10,WangTPC14}.

\subsection{A finite selector theorem}

There is nevertheless a rigorous existence mechanism when the shift bundle
is fixed and finite.  It is useful to state it independently of the primes.

\begin{theorem}[Finite selector]
\label{thm:finite-coactivation-selector}
Let $K\ge2$ be fixed and let
$x(n)=(x_1(n),\ldots,x_K(n))\in\{0,1\}^K$.  Define
\[
 E(N):=\sum_{n\le N}\one_{\{\sum_i x_i(n)\ge2\}},
 \qquad
 C_{ij}(N):=\sum_{n\le N}x_i(n)x_j(n)
 \quad (i<j).
\]
Then, for every $N$,
\begin{equation}
 E(N)\le\sum_{i<j}C_{ij}(N),
 \qquad
 \max_{i<j}C_{ij}(N)\ge\frac{E(N)}{\binom K2}.
 \label{eq:finite-selector-bound}
\end{equation}
If $E(N)\to\infty$, there are a fixed pair $(i_*,j_*)$ and a sequence
$N_r\to\infty$ for which
\[
 C_{i_*j_*}(N_r)\ge \frac{E(N_r)}{\binom K2}.
\]
In particular, $C_{i_*j_*}(N)\to\infty$.
\end{theorem}

\begin{proof}
For each $n$,
\[
 \one_{\{\sum_i x_i(n)\ge2\}}
 \le \sum_{i<j}x_i(n)x_j(n).
\]
Summing gives the first inequality in
\eqref{eq:finite-selector-bound}, and averaging over the $\binom K2$ pairs
gives the second.  Choose a maximizing pair for each $N$.  One pair occurs
for infinitely many $N$ because the set of pairs is finite.  Passing to
that subsequence proves the quantitative assertion.  Since $E(N_r)\to
\infty$ and each $C_{ij}(N)$ is nondecreasing, the last conclusion follows.
\end{proof}

In symbolic-dynamical language, the set of states with at least two active
coordinates is the finite union
\[
 \left\{x:\sum_i x_i\ge2\right\}
   =\bigcup_{i<j}\{x:x_i=x_j=1\}.
\]
Infinite recurrence to this union therefore forces infinite recurrence to
at least one fixed two-coordinate cylinder.  Finiteness is essential: the
argument does not invert a normalized average over a growing family of
shifts.

\begin{corollary}[Standard bounded-gap consequence at an unidentified shift]
\label{cor:unidentified-bounded-gap}
There exists a fixed even integer $h_*$ with
\[
 2\le h_*\le246
\]
such that $p$ and $p+h_*$ are simultaneously prime for infinitely many
primes $p$.
\end{corollary}

\begin{proof}
This is the Polymath8b bounded-gap theorem followed by finite pigeonhole;
it is not a new bounded-gap theorem and is not used in the analytic
reduction above.  Specifically, Theorems 16(i) and 17 of
\citet{Polymath2014} give
$\mathrm{DHL}[50;2]$ and $H(50)=246$.  Fix an admissible $50$-tuple
\[
 \cH=\{h_1<\cdots<h_{50}\}\subset[0,246]
\]
of diameter $246$, and apply
\cref{thm:finite-coactivation-selector} to
$x_i(n)=\one_{\mathbb P}(n+h_i)$.  The $\mathrm{DHL}[50;2]$ assertion says
that $E(N)\to\infty$, so one fixed pair $(i_*,j_*)$ is coactive infinitely
often.  Take $h_*=h_{j_*}-h_{i_*}$.  Admissibility modulo $2$ forces all
$h_i$ to have the same parity; hence $h_*$ is even, and its diameter gives
$2\le h_*\le246$.
\end{proof}

The corollary is an exact prime--prime recurrence statement, but its
quantifier is ``there exists $h_*$''.  It neither identifies $h_*$ nor
permits the substitution $h_*=2$.

\subsection{Why the aggregate cannot identify the coordinate}

The loss of the shift label is already exact in the abstract finite-channel
model.

\begin{proposition}[Aggregate non-identifiability]
\label{prop:aggregate-nonidentifiability}
Let $K\ge3$.  Knowledge of the entire aggregate path
$\{E(N):N\ge1\}$, even together with
$\{\sum_{i<j}C_{ij}(N):N\ge1\}$, does not determine which pair is
coactive infinitely often on the class of binary $K$-channel sequences.
\end{proposition}

\begin{proof}
Fix an infinite set $S\subset\N$ and two distinct coordinate pairs
$e\ne f$.  In the first sequence set precisely the two coordinates in $e$
equal to $1$ when $n\in S$, and set every coordinate to $0$ otherwise.
Construct the second sequence in the same way with $f$ in place of $e$.
For both sequences and every $N$,
\[
 E(N)=\sum_{i<j}C_{ij}(N)=|S\cap[1,N]|,
\]
but their unique recurrent coactivation pairs are different.
\end{proof}

This proposition is an information statement about the aggregate
observable, not a claim of logical independence for the actual primes.
Recovering a prescribed shift requires label-sensitive input that is absent
from the aggregate.

\subsection{The unresolved packet}

The Ramanujan calibration and Vaughan decomposition above reduce the
endpoint error to the displayed balanced Type-II packet, up to terms already
shown to be $o(X)$.  Recall that
$\cB_{h,\delta}=-\cH^{\rm HB}_{h,Q;U,V}$.  The present paper does not prove
the required estimate
\begin{equation}
 \cH^{\rm HB}_{h,Q;U,V}(X)=o(X)
 \label{eq:hard-packet-unproved}
\end{equation}
for that packet.  Consequently it does not prove
\eqref{eq:hl-endpoint-boundary}, a positive lower bound for
$\mathcal P_2(X)$, or the twin-prime conjecture.  The decomposition is a precise
information gate: Type-I distribution and the finite Ramanujan spectrum
have been discharged, while genuinely balanced, label-sensitive bilinear
control remains.  Known parity-barrier analyses explain why such additional
information cannot be replaced by formal sieve rearrangement alone
\citep{Selberg1952Parity,FordMaynard2024}.

For clarity, the final ledger is therefore:
\begin{center}
\begin{tabular}{@{}>{\raggedright\arraybackslash}p{0.29\textwidth}
                    >{\raggedright\arraybackslash}p{0.18\textwidth}
                    >{\raggedright\arraybackslash}p{0.42\textwidth}@{}}
\toprule
Object & Natural scale & Status here \\
\midrule
Complete row-centered residual & $X\log X$ & Explicit asymptotic; still
summed over all primitive orientations. \\
Fixed-shift prime endpoint & $X$ & Upper-bound scale only; no
Hardy--Littlewood asymptotic. \\
Finite admissible shift bundle & recurrence & Some fixed even
$h_*\le246$ is selected; its identity is not recovered. \\
Balanced hard packet & $X$ target & Isolated exactly; the necessary
$o(X)$ estimate is not proved. \\
\bottomrule
\end{tabular}
\end{center}
